\documentclass{IOS-Book-Article}

\usepackage[T1]{fontenc}
\usepackage{mathptmx}
% Latin Modern typewriter: Type 1 outlines in T1 (the default T1 cmtt comes out as Type 3 bitmaps)
\renewcommand{\ttdefault}{lmtt}
\usepackage{soul}\setuldepth{article}
\usepackage{longtable}
\usepackage{array}
\def\hb{\hbox to 11.5 cm{}}
\urlstyle{same}
\hypersetup{pdftitle={LexSR: Consolidating Legislation and Linking Case Law to Build a Legal Knowledge Graph for Suriname, a Low-Resource Jurisdiction},
  pdfauthor={Ernst van Gassen}}
\def\UrlBreaks{\do\/\do\-}
\newcommand{\corpusurl}{\url{https://doi.org/10.5281/zenodo.23467651}}

\begin{document}

%\pagestyle{headings}
\def\thepage{}
\begin{frontmatter}

\title{LexSR: Consolidating Legislation and Linking Case Law to Build a Legal Knowledge Graph for Suriname, a Low-Resource Jurisdiction}

%\markboth{}{2026\hb}

\author[A]{\fnms{Ernst} \snm{van Gassen}%
\thanks{Corresponding Author: Ernst van Gassen, E-mail: evg.gassen@gmail.com.}}

%\runningauthor{E. van Gassen}
%\runningtitle{LexSR: A Legal Knowledge Graph for Suriname}
\address[A]{Independent researcher at Arktos AI Lab}

\begin{abstract}
Access to consolidated, machine-readable legislation and linked case law is a prerequisite
for effective legal information systems. No such resource previously existed for Suriname, a
Dutch-speaking jurisdiction of about 640,000 people whose legal archive is mostly scanned.
We present LexSR, a pipeline that ingests raw legislative PDFs and
the Hof van Justitie's case-law portal, applies OCR, classifies documents, automatically
consolidates \emph{wijzigingswetten} (amendment acts) into dated versions of each act,
modelled on the Work and Expression levels of FRBR (Functional Requirements for
Bibliographic Records), and links case law to the legislation it applies. The pipeline produces a
corpus of 1,585 works (1,285 legislative, 300 judicial) and 68,127 fragments, from
which we construct a unified knowledge graph of 2,653 edges spanning amendment,
citation, and case-application relationships. We report two evaluations: a synthetic
citation-resolution benchmark, where our resolver passes 62 of 65 curated cases at article
level (100\% precision, 93.5\% recall), and, as the paper's methodological contribution, a measurement of
how much of that accuracy holds on the real, messily titled corpus the resolver was built
for. Only 21\% of extracted citations resolve to an
in-corpus work, against 11\% for exact title matching. We attribute the gap to a structural
mismatch between the canonical identifiers assumed by curated benchmarks and the
source-derived titles of scraped, OCR'd documents, and to gazette references, which the
corpus does not yet model. Amendment parsing yields a classified instruction for 25 of 36
amendment documents (69\%); most failures trace to operative text missing from the
ingested source. We discuss
implications for low-resource legal NLP and make the code, evaluation scripts and the
government-sourced part of the corpus publicly available.
\end{abstract}

\begin{keyword}
legal knowledge graph\sep legislative consolidation\sep case-law citation linking\sep
low-resource legal NLP\sep FRBR\sep Suriname
\end{keyword}
\end{frontmatter}
\pagestyle{empty}
\thispagestyle{empty}

\section{Introduction}

Consolidation (merging an amendment act's instructions into the base law it modifies) and
citation linkage (connecting case law to the provisions it applies) are both legally
critical and technically demanding. Most European jurisdictions maintain authoritative
consolidated legislation and case-law databases through governmental portals (e.g.,
\emph{wetten.overheid.nl} and \emph{rechtspraak.nl} in the Netherlands, EUR-Lex for EU law). Small and
developing jurisdictions often lack the institutional capacity for either, leaving
practitioners to track amendment chains and cross-reference case law by hand, or to rely
on commercial services, with no public integrated view of the law in force.

Suriname, a Dutch-speaking jurisdiction of about 640,000 people, shows this gap clearly. The \emph{Staatsblad} (official gazette; before
independence in 1975 the \emph{Gouvernementsblad}) is available as
individual scanned PDFs, scattered over government sites and a private legal portal, and the Hof van Justitie
publishes case decisions on a dedicated portal (\emph{rechtspraak.sr}, launched 2019), but no
system links the two: a decision citing ``artikel 4 Arbeidswet 1963'' offers no way to jump to
that provision, and no portal shows any act's consolidated text as amended.

This paper makes four contributions:

\begin{enumerate}
\item \textbf{LexSR corpus}: a structured, FRBR-modelled database of 1,285 Surinamese legislative
   works and 300 Hof van Justitie decisions, with 68,127 fragments, derived from eleven
   government sources, the private legal portal SRiS, the judicial case-law portal, and
   locally digitised archives.
\item \textbf{Amendment parser}: a hybrid regex + LLM system that extracts structured consolidation
   instructions from \emph{wijzigingswetten}, evaluated on every amendment act in the corpus.
\item \textbf{Case-law integration}: a scraper and ingestion pipeline for \emph{rechtspraak.sr} that
   segments decisions into numbered considerations and links them to cited legislation,
   reusing the same citation-resolution machinery built for legislation-to-legislation
   citations.
\item \textbf{Knowledge graph and an accuracy audit}: a unified graph of amendment, citation,
   and case-application edges over the full corpus, together with a measurement of the gap
   between synthetic-benchmark citation accuracy and real-corpus linkage rate, and an
   account of why the two diverge.
\end{enumerate}

Code, evaluation scripts and the government-sourced part of the corpus are
public.\footnote{Code and evaluation scripts:
\url{https://github.com/arktos-ai-lab/lexsr-suriname-legal-corpus}. The works obtained from
government sources (the 300 decisions and 751 legislative works) are archived with their
source files under CC BY 4.0 at \corpusurl; works obtained from SRiS are not redistributed.}

\section{Related Work}
\label{sec:related}

\subsection{Legal Document Modelling}

The FRBR (Functional Requirements for Bibliographic Records) model provides a
well-established abstraction for legislative documents, distinguishing the abstract Work
(a law), its Expression (a version at a given date), its Manifestation (a particular
encoding), and individual Items \cite{ifla1998}. Akoma Ntoso \cite{palmirani2011a} is
the primary XML schema for legislation structured around these concepts and has been adopted
by the UN, EU, and several national parliaments. The European Legislation Identifier (ELI)
\cite{pubsoffice2012} provides URI conventions compatible with this model. Where a provision
sits in this model varies: FRBRoo treats a part of a text as an Expression Fragment
\cite{frbroo}, while component-level approaches make each provision a Work in its own right
\cite{demartim2025}. We follow the first reading and extend it to case law: a court decision
is a Work with a single Expression, and its numbered considerations are fragments, the
case-law analogue of legislative articles.

\subsection{Legal Knowledge Systems for Consolidation}

Automated consolidation has received sustained attention in the legal knowledge systems
community. Biagioli, Francesconi et al. \cite{francesconi2004} describe the \emph{Norme in
Rete} project's XML standards and tools for cross-referencing and consolidating Italian
legislation. Spinosa et al. \cite{spinosa2009} apply NLP-based metadata extraction to detect
and classify textual amendments (repeal, substitution, integration) for the same
consolidation task. The \emph{legislation.gov.uk} system \cite{sheridan2010}
handles consolidation for UK law at scale with a semi-automated editorial workflow. Vitali
and Zeni \cite{vitali2007} discuss amendment formalisation in Akoma Ntoso. The Dutch
\emph{wetten.overheid.nl} portal maintains consolidated texts for all Dutch and former-colonial
legislation through a government editorial process, but does not publish its consolidation
tooling \cite{koop2022}. LexSR differs from these in targeting a low-resource setting (no
existing structured archive, predominantly scanned PDFs, and no editorial support) with a
fully automated pipeline.

\subsection{Legal Knowledge Graphs and Citation Networks}

Citation network analysis of case law has a long history in AI and Law, from early work on
precedent networks to more recent graph-based approaches to legal information retrieval and
argumentation mining \cite{benchcapon2012}. Legal knowledge graphs unifying legislation
and case law have been built for well-resourced jurisdictions; Dhani et al.
\cite{dhani2021}, for example, construct one from Indian court judgments and statutes for
case-similarity recommendation, and Lex Graph \cite{lexgraph} builds one over UK
legislation. Linkoln \cite{bacci2019} detects and links legislative
references in Italian texts and is deployed on the Italian Senate website. These approaches
resolve references against structured sources with an authoritative naming scheme; LexSR
must first create the identifiers it resolves to, from scanned PDFs and an unindexed case
portal.

\subsection{NLP for Legal Amendment Text}

Information extraction from legal amendment provisions has been studied for several European
languages. Lesmo, Mazzei, Palmirani et al. \cite{tulsi2013} present TULSI, an NLP system
combining deep syntactic analysis with a shallow semantic interpreter to extract and classify
modificatory provisions in Italian legislative texts. For Dutch, Boer et al. \cite{boer2007}
propose the MetaLex/LKIF XML standards for structuring legislation, and Hoekstra et al.
\cite{hoekstra2007} extend this with the LKIF Core ontology of basic legal concepts, but
explicit amendment/consolidation parsing for Dutch has received limited attention in the
JURIX and ICAIL literature. The use of large language models for
legal text understanding is by now standard \cite{chalkidis2020,koreeda2021}. Our hybrid
approach, regex-first with an LLM fallback for low-confidence cases, is motivated by the need for determinism and auditability while retaining flexibility for novel
constructions.

\subsection{Legal Corpora for Small Jurisdictions}

Multilingual legal corpora remain concentrated on the major European languages
\cite{niklaus2024}, and Krasadakis et al. \cite{krasadakis2024} survey the challenges
this poses for low-resource legal NLP. For Suriname, no machine-readable corpus linking
legislation and case law had been published before LexSR.

\section{Corpus and Ingestion Pipeline}
\label{sec:corpus}

\subsection{Sources}

LexSR ingests legislation from eleven government sites (parliament DNA.sr, the tax and
maritime authorities, the standards bureau, the government portal, six ministries), the
private legal portal SRiS (Stichting voor de Rechtsorde in Suriname), and locally digitised
PDF archives \cite{sources}. Of the 1,285 legislative works, 704 take their text from PDFs
published by parliament, 47 from other government sites, 499 from SRiS and 35 from local
archives. For the 215 acts available from both parliament and SRiS the parliament copy is
used, as SRiS, a private foundation, publishes copies, not authenticated texts. Six further
sites our scrapers support were excluded because their configured URLs no longer match the
live site structure (Section~\ref{sec:limitations}). Case law is sourced separately from
\emph{rechtspraak.sr}, the Hof van Justitie's decision portal, described in Section~\ref{sec:caselaw-scraping}. Legislative
source files are retained unchanged and tied to their works through the filename-derived
identifier; case-law works record source URL and retrieval time. A document obtained under
several file names is kept once: works with identical text, or with the same gazette issue
and near-identical OCR text, are merged. A SHA-256 manifest of
all source files is in the repository \cite{repo}. Article 11 of the Surinamese Auteurswet 1913
excludes legislation, judicial decisions and administrative decisions from copyright;
documents outside that scope (theses, standards drafts, third-party technical codes)
are released without their text.

\subsection{Data Model}

The canonical store maps FRBR entities as shown in Table~\ref{tab:frbr}. Items are not
modelled, and each fragment belongs to one Expression. The schema defines Manifestations,
but the corpus does not yet populate them.

\begin{table}[h]
\caption{FRBR entities and LexSR entities}
\label{tab:frbr}
\begin{tabular}{|l|l|l|l|}
\hline
\textbf{FRBR entity} & \textbf{LexSR entity} & \textbf{Legislation example} & \textbf{Case-law example} \\
\hline
Work & \texttt{Work} & Arbeidswet 1963 & SRU-HvJ-2022-39 \\
Expression & \texttt{Expression} & version\_date=2024-01-15 & single expression \\
Manifestation & \texttt{Manifestation} & source PDF & source HTML page \\
Expression Fragment & \texttt{Fragment} & Artikel 4, lid 2 & consideration 2.1 \\
\hline
\end{tabular}
\end{table}

Stable URIs follow an ELI-inspired pattern rooted at \texttt{/id/\{scheme\}/\{id\}}, with
\texttt{law} and \texttt{caselaw} as the two schemes currently populated. The amendment relation is
represented by an \texttt{AmendmentLink} table connecting an amending Work to the base Work it
modifies; the citation relation (legislation-to-legislation, case-to-legislation, and
case-to-case) is represented by a single generic \texttt{Citation} table typed by
\texttt{citation\_type}, avoiding a schema split between the two document kinds.

\subsection{Corpus Statistics}

Table~\ref{tab:stats} summarizes the corpus. The legislative works whose gazette issue is known (991) span 1904 to 2026; two thirds
date from 2010 onward, when digital-native PDFs became available. The case-law
sample spans 1952 to 2025 and was drawn from a discoverable pool of 900+ decisions on
\emph{rechtspraak.sr} (Section~\ref{sec:caselaw-scraping}); this paper reports on the first 300 ingested.

\begin{table}[h]
\caption{Corpus statistics}
\label{tab:stats}
\begin{tabular}{|l|r|}
\hline
\textbf{Metric} & \textbf{Value} \\
\hline
Legislative works & 1,285 \\
Case-law works (decisions) & 300 \\
Expressions & 1,594 \\
Fragments (articles, paragraphs, considerations) & 68,127 \\
Considerations extracted from case law & 2,684 \\
Citations extracted & 12,343 \\
Citations resolved to an in-corpus work & 2,613 (21.2\%) \\
Amendment links identified & 40 \\
Amendment links fully applied by the engine & 9 \\
Knowledge graph nodes / edges & 1,585 / 2,653 \\
\hline
\end{tabular}
\end{table}

\subsection{OCR and Structure Extraction}

For scanned PDFs, text is extracted using a pipeline that applies standard PDF text
extraction first and falls back to Tesseract OCR \cite{smith2007} for image-based pages,
after turning sideways pages upright and splitting scans that hold two gazette pages. Article boundaries
are detected heuristically using the standard Dutch/Surinamese legislative structure:
``Artikel N'' headings, ``Lid N'' sub-paragraphs, and lettered or numbered sub-items. OCR quality
varies substantially across the corpus, and a subset carries the lower text quality of an
English rather than a Dutch recognition model. We record this as a property of the
corpus, since it bears on the parsing coverage reported in Section~\ref{sec:parsing-eval}.

\section{Amendment Consolidation}
\label{sec:consolidation}

\subsection{Structure of a \emph{Wijzigingswet}}

A Surinamese amendment act (\emph{wijzigingswet} or \emph{wijzigingsbesluit}) follows a standard
template. The operative article names the base law and lists labelled instructions, as in
this excerpt from S.B.~2002 no.~113 \cite{sb2002113}:

\begin{quote}
\small\ttfamily
ARTIKEL I\\
In de Wet Huurwaardebelasting 1995 (S.B. 1995 no.108) worden\\
de volgende wijzigingen aangebracht:\\[4pt]
A. Artikel 5 lid 2 wordt vervangen door: De belasting wordt\\
\hspace*{3ex}ten minste vastgesteld op Sf. 20.000,-.\\
B. In artikel 6 lid 1 wordt de zinsnede `de vrije\\
\hspace*{3ex}verkoopwaarde [...]' vervangen door: [...]\\
C. Artikel 7 lid 1 komt te luiden als volgt: [...]
\end{quote}

Later articles (Artikel II, ...) may amend a second base law or
specify the in-force date. The act typically concludes with a \emph{Memorie van Toelichting}
(explanatory memorandum), which must be excluded from instruction extraction.

\subsection{Parser Architecture}

The \texttt{AmendmentParser} operates in two passes. \textbf{Pass 1 (regex)} splits the document into
sections by ``Artikel N'' boundaries (after stripping explanatory-memorandum text), extracts
the target law name from the section header or preamble, splits labelled instruction blocks
(A., B., C. or Roman-numeral labels), and classifies each block against a priority-ordered
set of $\sim$20 regex patterns covering the standard Dutch amendment constructions (repeal,
replacement, insertion, phrase substitution). Confidence is a fixed prior per pattern, not a
learned score: 0.95 for an explicit repeal, 0.70 to 0.90 for substitutions and insertions
whose new text was isolated, 0.45 to 0.55 when the pattern matched but the new text could
not be isolated, and 0.20 for an unclassified block. \textbf{Pass 2 (LLM fallback)}, used
only when an instruction's confidence falls below 0.70, re-sends the section to an LLM
with a structured JSON prompt (in the repository \cite{repo}); LLM output overrides regex output where its confidence is
higher. The fallback is disabled for every result reported here, so all figures below
describe regex-only behaviour and the fallback itself is not evaluated.

\subsection{Consolidation Engine}

\texttt{ConsolidationEngine} consumes \texttt{AmendmentInstruction} objects and applies them to the current
expression of the base law to produce a new \texttt{Expression} at the amendment's effective date.
Base-law resolution maps the amendment's target-law reference to a Work by title matching,
since the stated law name rarely matches the filename-derived work ID.
Let $n(s)$ lowercase $s$, delete the Dutch articles and prepositions \emph{de, het, van,
der, den, een} and collapse whitespace. A reference $r$ resolves to a work with title $t$ by
the first step that succeeds: (i) $r = t$ ignoring case; (ii) $n(r) = n(t)$. Steps (iii) and
(iv) need a law name in $r$ once article, \emph{lid} and \emph{sub} pointers are removed, and
skip amending acts (titles with \emph{wijziging}, \emph{verbeterblad}, \emph{intrekking})
unless $r$ names one: (iii) one of $n(r), n(t)$ contains the other, both at least four
characters, the closest in length winning; (iv) $r$ and $t$ share a four-digit year and $t$
contains a word of $n(r)$ longer than three characters. This
matching logic is shared with citation resolution, which requires the identical fallback for
the identical reason (Section~\ref{sec:linkage-gap}).

\section{Case-Law Integration and Knowledge Graph}
\label{sec:caselaw}

\subsection{Scraping \emph{rechtspraak.sr}}
\label{sec:caselaw-scraping}

\emph{rechtspraak.sr} is a WordPress-based site with no listing or search API; the only reliable
enumeration mechanism is the WordPress XML sitemap, which we parse recursively (sitemap
index $\rightarrow$ per-type sub-sitemaps $\rightarrow$ individual decision URLs). This alone surfaced 900+
decisions spanning 1952 to 2025. Each decision page exposes clean structured metadata (court,
case number, decision and publication dates, legal area) as label/value pairs, and full
decision text as HTML, considerably cleaner input than the OCR'd legislative PDFs, since
no scanning or recognition step is involved.

\subsection{Consideration Segmentation}

Surinamese court decisions carry no explicit ``r.o.'' (\emph{rechtsoverweging}) label, the
marker of each numbered consideration in Dutch judgments. Recent
decisions number their reasoning under a ``De beoordeling'' (assessment) heading, as
section.paragraph (``2.1'', ``4.7'') or as plain paragraphs (``1.'', ``2.''); older decisions
give it under ``Ten aanzien van het recht'' as paragraphs opening with ``Overwegende, dat''.
We segment on these markers up to the dictum; the text before the reasoning and the dictum
are kept as two further fragments, so no text is dropped. This splits 286 of the 300
decisions into 2,670 considerations; the other 14 fit neither layout and are kept whole
(2,684 consideration fragments in all).

\subsection{Citation Extraction from Case Law}

Legislative citations in case-law text (``artikel 4 lid 2 Arbeidswet 1963'') are extracted
with the same resolver and pipeline as legislation-to-legislation citations, with no
case-law-specific logic, yielding 583 case-to-legislation citations across the 300
decisions. Case-to-case citations by the portal's own identifiers (``SRU-HvJ-2019-3'') use a
separate regex; none were found, because decisions cite each other by docket number
instead (Section~\ref{sec:limitations}).

\subsection{Knowledge Graph Construction}

We assemble a unified graph directly from the canonical store, adding no schema or
storage layer, with one node per Work (typed \texttt{legislation} or \texttt{case\_law}) and three edge
types: AMENDS (from \texttt{AmendmentLink}), CITES (legislation-to-legislation citations), and
APPLIES (case-to-legislation citations). The resulting graph over the full corpus has
1,585 nodes and 2,653 edges (2,178 CITES, 435 APPLIES, 40 AMENDS), with a largest connected
component of 687 nodes. The ten most-cited works are, in order: Personeelswet (civil service
law), Burgerlijk Wetboek (Civil Code), Wetboek van Strafrecht (Criminal Code), Grondwet
(Constitution), Kiesregeling (electoral law), Wetboek van Strafvordering (Code of Criminal
Procedure), Wetboek van Koophandel (Commercial Code), Wet Loonbelasting (wage tax),
Inkomstenbelasting 1922 (income tax) and Wet Regionale Organen (regional bodies). This distribution matches expectation for a general legal system and offers a
face-validity check on the extraction pipeline independent of the benchmark in Section~\ref{sec:resolver-eval}.

\section{Evaluation}
\label{sec:eval}

\subsection{Citation Resolution: Synthetic Benchmark}
\label{sec:resolver-eval}

We evaluate the citation resolver on a 65-case known-answer dataset (61 citation and
4 temporal cases, including gazette references and deliberately unresolvable negatives),
whose expected answers we checked by hand (all correct). Most cases were generated with an AI coding assistant together
with the resolver patterns; six are quoted from decisions in the corpus, and decisions are
cited by their \emph{rechtspraak.sr} identifiers, as Suriname has no ECLI. Cases naming an
article are scored at article level, ignoring the version date (the temporal cases thus do
not test the version). The resolver passes 62 of 65 (precision 100\%, recall 93.5\%, F1
96.7\%). The three misses are a Dutch Hoge Raad judgment cited by ECLI (Surinamese
decisions also cite Dutch case law), where the resolver keeps the sentence-final full stop
in the identifier; an article separated from its law by a bracketed gazette reference; and
the Wetboek van Burgerlijke Rechtsvordering, which the abbreviation registry lacks. Because patterns and most cases share an
origin, this result
shows pattern coverage on well-formed input, not generalisation; the next section removes
that condition.

\subsection{Citation Resolution: Real-Corpus Linkage Gap}
\label{sec:linkage-gap}

Synthetic-benchmark accuracy measures whether the resolver identifies \emph{that} a
reference exists and which canonical URI it maps to. It does not measure whether that URI
names anything in the corpus. The two come apart because the benchmark and the corpus use
different names: the resolver's abbreviation registry maps ``Grondwet'' to the canonical
slug \texttt{grondwet-1987}, while a scraped corpus with no naming authority identifies its
works by their source documents (\texttt{grondwet-met-mvt-\allowbreak{}4fb9cced}, titled
``Grondwet met Mvt''). No registry slug is the identifier of a work in the corpus, so
canonical lookup alone links nothing. Resolution therefore falls back on the title matching
of Section~\ref{sec:consolidation}, which prefers the closest title and skips amending acts.

With this fallback, 2,613 of the 12,343 citations in the corpus (21.2\%)
resolve to an in-corpus work. As a baseline, linking each extracted
reference to a work by exact, case-insensitive title match links 1,363 citations (11.0\%).
Checked by hand on a sample, 96.5\% of the links are correct (40 of 48 distinct law-to-work pairs), as are all 14 decidable pairs that link to a DNA.sr text; the
errors link to an implementing decree, match on the year alone, or read a tax name as a law
reference. The remaining 79\% are overwhelmingly gazette
references (\emph{Staatsblad}/\emph{Gouvernementsblad} citations). The resolver does not
link these, because the corpus does not model gazette issues as entities. Fuzzy title
matching is no substitute: its substring step pairs ``G.B.~1973 no.~9'' with an act numbered
1973 no.~73. Yet 2,625 unresolved gazette references name a gazette, year and number that
appear together in the title of a work in the corpus; checked by hand, that work is the cited
act for all 74 of 75 sampled references that could be decided. An exact gazette-number match
is thus a direct next step (Section~\ref{sec:future}).

A resolver validated only against a curated benchmark can thus report near-perfect accuracy
while linking little of its corpus. We suggest a checklist item for similar projects:
\emph{measure resolution against the real corpus's actual identifiers, not just against
curated ground truth.}

\subsection{Amendment Parsing Coverage}
\label{sec:parsing-eval}

We evaluate the amendment parser (regex pass only) on all 36 amendment documents in the
corpus, per document rather than per \texttt{AmendmentLink} row (40 links, since a document
may target more than one base law). A document counts as parsed when the regex pass extracts
at least one classified instruction from it, that is, an instruction other than the 0.20
fallback; whether the extracted instructions are correct is not assessed by this measure.
25 of the 36 documents (69.4\%) are parsed.

Checked by hand, 9 of the 11 failing documents contain no operative instruction in the stored
text: it holds only the closing formula (in-force clause and signatures), the explanatory
memorandum, a bare enumeration without its amending sentence, or OCR too damaged to read.
One repeals an earlier decree in its closing provisions (``vervalt het Presidentieel
Besluit \ldots''), which the patterns miss, and in one the OCR misread the article numbers
(``Artikel | lid I'' for ``Artikel 1 lid 1''). Parsing failures thus mostly reflect what
ingestion captured from the source, not the parser rules.

\subsection{Engine Validation}

Nine amendments have been fully applied by the consolidation engine and are available for inspection in the canonical store; the remaining
parseable documents fail engine-level base-law resolution or produce low-confidence
instructions withheld pending review.

\section{Discussion and Future Work}
\label{sec:discussion}

\subsection{Limitations}
\label{sec:limitations}

\textbf{Source stability.} Six candidate sources were excluded because their URLs or site
structure changed; such sites rarely expose a stable API, so scrapers decay silently.

\textbf{Amendment parsing and upstream stages.} The 69\% coverage figure
(Section~\ref{sec:parsing-eval}) counts documents with at least one classified instruction
and says nothing about instruction correctness. OCR, document classification and
downstream tasks such as retrieval are not evaluated separately.

\textbf{Other jurisdictions.} The pipeline is untested outside Suriname. For Spanish-, French-
or Portuguese-language gazettes in Latin America or Africa we expect OCR to carry over and
scrapers and patterns to need rewriting; Indigenous-language texts, with few OCR models or
evaluation data, are harder.

\textbf{Case-to-case citation coverage.} The extractor found no case-to-case citations
(Section~\ref{sec:caselaw}) because it matches only the portal's own identifiers, yet 59 of
the 300 decisions cite a first-instance judgment by docket number (``A.R. no.~11-4959'').
Linking appeal chains needs this pattern and the first-instance decisions on the portal.

\subsection{Effort, Operation and Generalisability}

LexSR was built by one researcher with legal training (Dutch tax law) but no specialist
knowledge of Surinamese law, working with AI coding assistants; no Surinamese legal expert
reviewed the output. The scrapers, regex patterns, abbreviation registry and 65-case gold
set were produced this way, and corpus text was not manually corrected. New publications
are added by re-running the scrapers, which skip stored documents; we have not measured
throughput or compute cost. Porting LexSR to another
jurisdiction means replacing four jurisdiction-specific parts: scraper configurations, the
citation abbreviation registry, the gazette-reference patterns, and the $\sim$20 amendment
patterns. Suriname shares Dutch-language drafting templates with Aruba, Cura\c{c}ao and
Sint Maarten, where the last two should carry over largely unchanged.

\subsection{Future Work}
\label{sec:future}

Next steps are: replacing fuzzy title matching with a resolution service of the kind
URN:NIR provides for Italian law, mapping canonical names to stored documents through a
maintained alias table; taking titles and identifiers from the printed heading rather than
the file name; linking gazette references by exact year and number, then modelling
\emph{Staatsblad}/\emph{Gouvernementsblad} issues as entities; ingesting the remaining 600+
decisions on \emph{rechtspraak.sr}; an instruction-level gold standard for the amendment
parser; and automated detection of scraper configuration decay.

\section{Conclusion}
\label{sec:conclusion}

We presented LexSR, a versioned, consolidated corpus of Surinamese legislation and case law
with a knowledge graph over it. Its central contribution is methodological: we measured how
much of the resolver's synthetic-benchmark accuracy (100\% precision, 93.5\% recall) holds on
the real corpus (21.2\% linkage, against 11.0\% for exact titles), and traced the gap to the
mismatch between canonical and source-derived identifiers and to gazette references. Amendment parsing covers 25 of 36 amendment
documents; most failures trace to text that ingestion did not capture.

\begin{thebibliography}{99}

\bibitem{ifla1998}
IFLA Study Group. Functional Requirements for Bibliographic Records. IFLA Publications; 1998.

\bibitem{palmirani2011a}
Palmirani M, Vitali F. Akoma-Ntoso for Legal Documents. In: Sartor G, et al., editors.
Legislative XML for the Semantic Web. Springer; 2011. p. 75-100.

\bibitem{pubsoffice2012}
Publications Office of the European Union. European Legislation Identifier (ELI). Official
Journal 2012/C 325/02; 2012.

\bibitem{francesconi2004}
Biagioli C, Francesconi E, Spinosa P, Taddei M. XML Documents Within a Legal Domain:
Standards and Tools for the Italian Legislative Environment. In: Document Analysis Systems
VI (DAS 2004). Lecture Notes in Computer Science, vol 3163. Springer; 2004.

\bibitem{spinosa2009}
Spinosa P, Giardiello G, Cherubini M, Marchi S, Venturi G, Montemagni S. NLP-based metadata
extraction for legal text consolidation. In: Proceedings of the 12th International Conference
on Artificial Intelligence and Law (ICAIL 2009); 2009. p. 40-49.

\bibitem{sheridan2010}
Sheridan J, Tennison J. Linking UK government data. In: Proceedings of LDOW 2010; 2010.

\bibitem{vitali2007}
Vitali F, Zeni F. Towards a country-independent data format: the Akoma Ntoso experience. In:
Proceedings of the V Legislative XML Workshop. Florence; 2007.

\bibitem{koop2022}
KOOP. Offici\"{e}le bekendmakingen: wetten.overheid.nl. Kennis- en exploitatiecentrum voor
offici\"{e}le overheidspublicaties; 2022.

\bibitem{frbroo}
Bekiari C, Doerr M, Le Boeuf P, Riva P, editors. FRBR object-oriented definition and
mapping from FRBRER, FRAD and FRSAD, version 2.4. IFLA; 2015.

\bibitem{demartim2025}
de Martim H. Modeling the Diachronic Evolution of Legal Norms: An LRMoo-Based,
Component-Level, Event-Centric Approach to Legal Knowledge Graphs. arXiv:2506.07853; 2025.

\bibitem{lexgraph}
Incubator for Artificial Intelligence (i.AI), UK Government. Lex Graph: building a knowledge
graph of UK legislation [software]; 2025. Available from:
\url{https://github.com/i-dot-ai/lex-graph}.

\bibitem{bacci2019}
Bacci L, Agnoloni T, Marchetti C, Battistoni R. Improving Public Access to Legislation
Through Legal Citations Detection: The Linkoln Project at the Italian Senate. In: Knowledge
of the Law in the Big Data Age. Frontiers in Artificial Intelligence and Applications.
IOS Press; 2019. p. 149-158.

\bibitem{benchcapon2012}
Bench-Capon T, et al. A history of AI and Law in 50 papers. Artificial Intelligence and Law.
2012;20(3):215-319.

\bibitem{dhani2021}
Dhani JS, Bhatt R, Ganesan B, Sirohi P, Bhatnagar V. Similar Cases Recommendation using Legal
Knowledge Graphs. arXiv:2107.04771; 2021.

\bibitem{tulsi2013}
Lesmo L, Mazzei A, Palmirani M, Radicioni DP. TULSI: an NLP system for extracting legal
modificatory provisions. Artificial Intelligence and Law. 2013;21:139-172.

\bibitem{boer2007}
Boer A, Winkels R, Vitali F. Proposed XML standards for law: MetaLex and LKIF. In:
Proceedings of JURIX 2007. IOS Press; 2007.

\bibitem{hoekstra2007}
Hoekstra R, Breuker J, Di Bello M, Boer A. The LKIF core ontology of basic legal concepts. In:
Proceedings of the LOAIT Workshop. CEUR Workshop Proceedings, vol 321; 2007. p. 43-63.

\bibitem{chalkidis2020}
Chalkidis I, et al. LEGAL-BERT: The muppets straight out of law school. In: Findings of
EMNLP 2020; 2020.

\bibitem{koreeda2021}
Koreeda Y, Manning C. ContractNLI: A dataset for document-level natural language inference
for contracts. In: Findings of EMNLP 2021; 2021.

\bibitem{niklaus2024}
Niklaus J, et al. MultiLegalPile: A 689GB Multilingual Legal Corpus. In: Proceedings of ACL
2024 (Volume 1: Long Papers); 2024.

\bibitem{krasadakis2024}
Krasadakis P, Sakkopoulos E, Verykios VS. A Survey on Challenges and Advances in Natural
Language Processing with a Focus on Legal Informatics and Low-Resource Languages.
Electronics. 2024;13(3):648.

\bibitem{sources}
Sources of the corpus (accessed 2026): De Nationale Assembl\'{e}e,
\url{https://www.dna.sr/wetgeving/surinaamse-wetten/}; Belastingdienst Suriname,
\url{https://belastingdienst.sr}; Maritieme Autoriteit Suriname, \url{https://mas.sr};
Suriname Standards Bureau, \url{https://www.ssb.sr}; Government of Suriname and its
ministries, \url{https://gov.sr}; Stichting voor de Rechtsorde in Suriname,
\url{https://www.sris.sr}; Hof van Justitie, \url{https://rechtspraak.sr}.

\bibitem{smith2007}
Smith R. An Overview of the Tesseract OCR Engine. In: Proceedings of the Ninth
International Conference on Document Analysis and Recognition (ICDAR 2007). IEEE; 2007.
p. 629-633.

\bibitem{sb2002113}
Wet van 27 december 2002, houdende wijziging van de Wet Huurwaardebelasting 1995.
Staatsblad van de Republiek Suriname 2002 no. 113. Available from:
\url{https://www.dna.sr/media/5iooda3s/sb-2002-no-113.pdf}.

\bibitem{repo}
van Gassen E. LexSR: code, evaluation scripts and source manifest [software]; 2026.
Available from: \url{https://github.com/arktos-ai-lab/lexsr-suriname-legal-corpus}.

\end{thebibliography}
\end{document}
